% 太阳系外行星（综述）
% license CCBYSA3
% type Wiki

本文根据 CC-BY-SA 协议转载翻译自维基百科\href{https://en.wikipedia.org/wiki/Exoplanet}{相关文章}。

\begin{figure}[ht]
\centering
\includegraphics[width=8cm]{./figures/c4a53fef76fc3e3b.png}
\caption{由 W. M. 凯克天文台在七年间拍摄的 HR 8799 系统中的四颗系外行星影像，行星的运动由年度观测结果插值显示} \label{fig_Exopla_2}
\end{figure}
\begin{figure}[ht]
\centering
\includegraphics[width=8cm]{./figures/fa92c1fc35546edf.png}
\caption{开普勒-37 系统中系外行星的尺寸与水星、火星和地球的大小对比} \label{fig_Exopla_1}
\end{figure}
系外行星（exoplanet，或称太阳系外行星）是指位于太阳系之外的行星。首个得到确认的系外行星发现于 1992 年，环绕一颗脉冲星运行；而围绕主序星的首次确认发现发生在 1995 年。另有一颗行星最早于 1988 年被探测到，并在 2003 年得到确认。2016 年，人们认识到关于系外行星的最早可能证据甚至可以追溯到 1917 年。截止至 2025 年 12 月 4 日，已确认的系外行星共有 6,053 颗，分布在 4,510 个行星系统中，其中 1,022 个系统包含多于一颗行星$^{[1][2]}$。

探测系外行星的方法有多种，其中凌日测光法和多普勒光谱法发现的数量最多。但这些方法存在明显的观测偏差，更容易发现靠近母恒星的行星；因此，已探测到的系外行星中约有 85\% 位于潮汐锁定区之内$^{[3]}$。据估计，大约每 5 颗类太阳恒星中就有 1 颗拥有一颗位于宜居带内的“类地尺寸”行星$^{[a][b][c][4][5]}$。若假设银河系中约有 2,000 亿颗恒星$^{[d]}$，则可以推测银河系内可能存在约 110 亿颗潜在宜居的类地尺寸行星；若将数量众多的红矮星系统也纳入考虑，这一数字可上升至约 400 亿颗$^{[6]}$。

目前已知质量最小的系外行星是 Draugr，其质量约为月球的两倍。NASA 系外行星档案中列出的质量最大的系外行星是 HR 2562 b$^{[7][8][9]}$，其质量约为木星的 30 倍。然而，根据某些以氘核聚变为界限的行星定义$^{[10]}$，它的质量可能过大，更应被归类为褐矮星。已知系外行星的轨道周期跨度极大，从不足一小时（最靠近母恒星者）到数千年不等。有些系外行星距离母恒星极远，以至于难以判断它们是否仍与恒星保持引力束缚。

距离地球最近的系外行星位于约 4.2 光年（1.3 秒差距）之外，环绕着比邻星运行，比邻星是距离太阳最近的恒星$^{[11]}$。在另一极端，人们也发现了关于河外行星的证据，即位于其他星系中的系外行星$^{[12][13]}$。

系外行星的发现极大地激发了人们对地外生命搜索的兴趣，尤其关注那些位于恒星宜居带（有时称为“金发姑娘区”）内的行星，因为在那里行星表面可能存在液态水——这是我们所知生命的必要条件。然而，对行星宜居性的研究还需要综合考虑诸多其他因素，以评估一颗行星是否适合孕育生命$^{[14]}$。詹姆斯·韦布空间望远镜（JWST）与地基及其他空间望远镜协同工作，有望在系外行星的成分、环境条件及宜居性等方面提供更深入的认识$^{[15]}$。

流浪行星是指不隶属于任何行星系统的行星。这类天体通常被视为独立于行星之外的一个类别，尤其当它们是气体巨行星时，往往被归类为次级褐矮星$^{[16]}$。银河系中流浪行星的数量可能高达数十亿颗$^{[17][18]}$。
\subsection{定义}
\subsubsection{IAU（国际天文学联合会）}
国际天文学联合会（IAU）所使用的“行星”这一术语的官方定义仅适用于太阳系，因此并不直接适用于系外行星$^{[19][20]}$。IAU 系外行星工作组（Working Group on Extrasolar Planets）于 2001 年发布了一份立场声明，其中给出了“行星”的工作性定义，并于 2003 年进行了修订$^{[21]}$。当时，系外行星被定义为满足以下条件的天体：
\begin{itemize}
\item 真质量低于氘发生热核聚变的临界质量（对于太阳金属丰度的天体，目前计算值约为 13 个木星质量），且环绕恒星或恒星遗骸运行的天体，被定义为“行星”（不论其形成机制如何）。判定一个系外天体是否为行星所需的最小质量或尺寸，应与太阳系中采用的标准相同。
\item 真质量高于氘发生热核聚变临界质量的亚恒星天体，被定义为“褐矮星”，不论其形成方式或所在位置如何。
\item 在年轻疏散星团中发现的、质量低于氘热核聚变临界质量且不受恒星束缚而自由漂浮的天体，不被视为“行星”，而应归类为“次级褐矮星”（或其他更合适的名称）。
\end{itemize}
这一工作性定义随后在 2018 年 8 月由 IAU 的 F2 委员会（系外行星与太阳系委员会）进行了修订$^{[22][23]}$。目前，IAU 对系外行星的官方工作定义如下：
\begin{itemize}
\item 真质量低于氘发生热核聚变的临界质量（对于太阳金属丰度天体，目前计算值约为 13 个木星质量），并且环绕恒星、褐矮星或恒星遗骸运行，同时其与中心天体的质量比低于 L4/L5 拉格朗日点不稳定性条件（$M/M_{\text{central}} < \frac{2}{25+\sqrt{621}}$）的天体，被定义为“行星”（不论其形成机制如何）。
\item 判定一个系外天体是否为行星所需的最小质量或尺寸，应与太阳系中采用的标准相同。
\end{itemize}
\subsubsection{替代方案}
IAU 的工作性定义并非始终被采用。另一种替代观点认为，应当依据**形成机制**来区分行星与褐矮星。普遍认为，巨行星通过“核吸积”形成，这一过程有时可能产生质量超过氘聚变阈值的行星$^{[24][25][10]}$；此类大质量行星可能已经被观测到$^{[26]}$。褐矮星则类似恒星，是由气体云的直接引力坍缩形成的，而这种形成机制同样可以产生质量低于 13 个木星质量、甚至低至 1 个木星质量的天体$^{[27]}$。
在这一质量范围内、以数百或数千天文单位（AU）的宽分离轨道绕恒星运行、且恒星/天体质量比很大的对象，很可能是以褐矮星方式形成的；其大气成分预计更接近其宿主恒星，而非通过吸积形成的行星（后者通常富含较重元素）。截至 2014 年 4 月，大多数被直接成像的“行星”质量较大且轨道很宽，因此很可能代表了褐矮星形成机制的低质量端$^{[28]}$。有研究提出，质量高于 10 个木星质量的天体是通过引力不稳定形成的，不应被视为行星$^{[29]}$。

此外，“13 个木星质量”的分界线并不存在精确的物理意义。氘聚变在某些质量低于该阈值的天体中同样可能发生$^{[10]}$，且发生的氘聚变量在一定程度上取决于天体的化学组成$^{[30]}$。2011 年，《系外行星百科全书》收录了质量高达 25 个木星质量的天体，并指出：“在观测到的质量谱中，13 个木星质量附近并不存在特殊特征，这进一步强化了放弃该质量上限的选择”$^{[31]}$。截至 2016 年，基于对质量—密度关系的研究，该上限被提高到 60 个木星质量$^{[32][33]}$。
Exoplanet Data Explorer 收录了质量最高达 24 个木星质量的天体，并附注：“IAU 工作组所采用的 13 个木星质量区分标准在物理上缺乏动机（尤其是对具有岩质核心的行星而言），且在观测上由于 $\sin i$ 不确定性而存在问题”$^{[34]}$。NASA 系外行星档案库则收录质量（或最小质量）不超过 30 个木星质量的天体$^{[35]}$。另一个区分行星与褐矮星的判据（不同于氘聚变、形成过程或轨道位置）是：其核心压力是由库仑压力主导，还是由电子简并压主导，两者的分界大约在 5 个木星质量附近$^{[36][37]}$。
\subsubsection{确认}
在 NASA 的系外行星档案库中，一颗系外行星的确认标准是：要么“不同的观测技术揭示了只能由行星解释的特征”$^{[38]}$，要么通过分析性技术加以确认$^{[2]}$。
对于《系外行星百科全书》而言，“若某一天体在经认可的论文或专业会议上被明确无歧义地主张为行星，则被视为已确认”$^{[39]}$。
\subsection{命名法}
\begin{figure}[ht]
\centering
\includegraphics[width=6cm]{./figures/a30e59faab567d01.png}
\caption{系外行星 HIP 65426b 是围绕恒星 HIP 65426 发现的第一颗行星$^{[40]}$。} \label{fig_Exopla_3}
\end{figure}
系外行星的命名惯例是国际天文学联合会（IAU）所采用的多恒星系统命名体系的延伸。对于绕单一恒星运行的系外行星，其 IAU 指定名称由母恒星的指定名或专名加上一个小写字母构成$^{[41]}$。字母按照围绕母恒星发现行星的先后顺序依次分配，因此在一个行星系统中，最先发现的行星被标记为“b”（母恒星本身视为“a”），随后发现的行星依次使用后续字母。如果同一系统中的多颗行星在同一时间被发现，则距离恒星最近的行星获得下一个字母，其余行星按轨道尺度顺序依次命名。IAU 还制定了一套临时认可的标准，用于处理环双星行星的命名。少数系外行星已被 IAU 正式赋予专名，此外还存在其他命名体系$^{[citation\ needed]}$。
\subsection{探测历史}
\begin{figure}[ht]
\centering
\includegraphics[width=8cm]{./figures/d13a18693bd77d1a.png}
\caption{截至 2022 年的现有及未来系外行星任务的 NASA 示意图。} \label{fig_Exopla_4}
\end{figure}
几个世纪以来，科学家、哲学家以及科幻作家一直怀疑系外行星的存在，但始终无法确定它们是否真实存在、数量是否普遍，或是否与太阳系行星相似。19 世纪曾提出过多种探测主张，但均被天文学家否定（有待引文）。

1917 年，人们记录到了围绕范·马南 2（Van Maanen 2）运行的一颗可能的系外行星的最早证据，但直到 2016 年才被认识为系外行星的迹象。天文学家沃尔特·西德尼·亚当斯（Walter Sydney Adams）使用威尔逊山天文台的 60 英寸望远镜获得了该恒星的光谱，并将其解释为 F 型主序星的光谱。后来，在对白矮星异常成分的研究中，这一光谱被重新审视。如今认为，这种光谱可能是附近一颗系外行星在恒星引力作用下被粉碎，其残余尘埃随后落到恒星表面所造成的结果。

20 世纪中期还出现过多起其他发现主张，涉及天鹅座 61、拉兰德 21185 以及巴纳德星，但这些主张直到 20 世纪 70 年代中后期才被否定（见下文“被否定的发现”）。1988 年，又有一次被认为是系外行星的疑似科学探测。随后不久，1992 年，亚历山大·沃尔什昌（Aleksander Wolszczan）和戴尔·弗雷尔（Dale Frail）宣布在毫秒脉冲星 PSR B1257+12 周围发现了两颗质量与类地行星相当的行星，这是目前被普遍接受的首次系外行星发现。1995 年，天文学家确认了第一颗围绕主序星运行的系外行星——一颗巨行星，以四天的周期绕附近恒星 51 飞马座运行。此后，一些系外行星被望远镜直接成像，但绝大多数仍是通过凌日法和径向速度法等间接方法探测到的。
\begin{figure}[ht]
\centering
\includegraphics[width=8cm]{./figures/85a4bfec6f369e50.png}
\caption{截至 2024 年 9 月，每年探测到的系外行星数量 $^{[46]}$} \label{fig_Exopla_5}
\end{figure}
2018 年 2 月，研究人员利用**钱德拉 X 射线天文台**并结合一种称为**微引力透镜**的行星探测技术，在一个遥远的星系中发现了行星存在的证据。他们指出：“其中一些系外行星（相对而言）小到只有月球大小，而另一些则巨大如木星。与地球不同，大多数系外行星并未紧密束缚在恒星周围，而是在太空中游荡，或在恒星之间以较为松散的方式运行。我们可以估计，这个［遥远］星系中的行星数量超过一万亿。”$^{[47]}$
\subsubsection{早期猜想}
“我们宣告这片空间是无限的……其中存在着无数与我们自身世界同类的世界。”

—— 乔尔达诺·布鲁诺（1584）$^{[48]}$

16 世纪，意大利哲学家乔尔达诺·布鲁诺是哥白尼学说（即地球和其他行星绕太阳运行的日心说）的早期支持者之一。他提出，恒星本身与太阳相似，同样也应当伴随着行星$^{[49]}$。

18 世纪，艾萨克·牛顿在其《自然哲学的数学原理》（*Principia*）结尾的《总注》（*General Scholium*）中也提及了这种可能性。他将恒星系统与太阳系作类比，写道：“如果恒星是类似系统的中心，那么它们将按照相似的设计被构造，并同样服从于唯一的主宰。”$^{[50]}$

1938 年，D. Belorizky 论证了利用凌日光度测量法搜寻“系外木星”（exo-Jupiters）的可行性$^{[51]}$。

1952 年，在第一颗“热木星”被发现之前 40 多年，奥托·斯特鲁维（Otto Struve）指出，并不存在任何必然理由认为行星不能比太阳系中的情况更靠近其母恒星。他还提出，可以通过多普勒光谱法和凌日法来探测处于短周期轨道上的超级木星行星$^{[52]}$。b6
\subsubsection{被否定的发现主张}
自 19 世纪以来，关于系外行星探测的主张便屡有提出。其中最早的一些涉及双星系统 70 蛇夫座。1855 年，东印度公司马德拉斯天文台的威廉·斯蒂芬·雅各布报告称，其轨道异常“高度可能”表明该系统中存在一个“行星体”$^{[53]}$。19 世纪 90 年代，芝加哥大学和美国海军天文台的托马斯·J. J.·西（Thomas J. J. See）声称，这些轨道异常证明在 70 蛇夫座系统中，围绕其中一颗恒星存在一个轨道周期为 36 年的暗天体$^{[54]}$。然而，福里斯特·雷·莫尔顿（Forest Ray Moulton）随后发表论文证明，在这些轨道参数下，该三体系统将高度不稳定$^{[55]}$。

关于 61 天鹅座可能拥有行星系统的主张也曾多次出现。1942 年，斯普劳尔天文台的卡伊·斯特兰德（Kaj Strand）观测到 61 天鹅座 A 与 B 的轨道运动存在微小但系统性的变化，因而推测在 61 天鹅座 A 周围必然存在一个质量约为 16 个木星质量的第三天体$^{[56]}$。此后又有多项类似主张，但较新的观测至今仍未给出确认。详见：61 Cygni：Claims of a planetary system。在研究 61 天鹅座的同时，天文学家也对拉朗德21185提出了类似的系外行星存在主张，见：Lalande 21185#Past claims of planets。

20 世纪 50 至 60 年代，斯沃斯莫尔学院的彼得·范·德·坎普（Peter van de Kamp）又发表了一系列颇具影响力的探测主张，认为巴纳德星周围存在行星$^{[57]}$。如今，天文学界普遍认为这些早期探测报告均为误判$^{[58]}$。

1991 年，安德鲁·莱恩（Andrew Lyne）、M. Bailes 和 S. L. Shemar 利用脉冲星计时变化，声称在PSR 1829−10周围发现了一颗脉冲星行星$^{[59]}$。这一消息一度引起广泛关注，但莱恩及其团队随后很快撤回了这一结论$^{[60]}$。
\subsubsection{已确认的发现}
\begin{figure}[ht]
\centering
\includegraphics[width=6cm]{./figures/7817ae950047ded9.png}
\caption{} \label{fig_Exopla_6}
\end{figure}